\section{Aplikasi Strategi \textit{Greedy}}
\label{sec:bab3}

\subsection{Pemetaan Persoalan Battlecode ke Elemen-Elemen Algoritma
            \textit{Greedy}}

Permainan Battlecode 2025 dapat dimodelkan sebagai sebuah persoalan
optimasi yang diselesaikan secara inkremental, di mana setiap robot
membuat keputusan aksi terbaik pada setiap ronde berdasarkan informasi
lokal yang tersedia.  Pemetaan antara elemen-elemen persoalan Battlecode
dan elemen-elemen algoritma \textit{greedy} adalah
sebagai berikut.

\subsubsection{Himpunan Kandidat}

Himpunan kandidat $C$ terdiri atas seluruh aksi yang secara teknis dapat
dilaksanakan oleh sebuah robot pada satu ronde.  Aksi-aksi tersebut
meliputi pilihan arah gerak ke delapan tile tetangga, pilihan tile yang
dapat diserang atau diwarnai dalam radius aksi, serta aksi-aksi khusus
sesuai kemampuan tipe robot masing-masing.

\subsubsection{Himpunan Solusi}

Himpunan solusi $S$ adalah rangkaian aksi yang telah dipilih dan
dilaksanakan oleh sebuah robot sepanjang permainan.  Himpunan ini
dibangun secara bertahap: setiap ronde, satu aksi baru ditambahkan ke
dalam $S$ berdasarkan hasil evaluasi fungsi seleksi.

\subsubsection{Fungsi Seleksi}

Fungsi seleksi menentukan aksi terbaik dari himpunan kandidat $C$ pada
setiap langkah.  Fungsi ini diwujudkan melalui heuristik penilaian yang
disesuaikan dengan tujuan masing-masing tipe robot: robot penyerang
memilih tile yang paling menguntungkan untuk diwarnai, robot pembersih
memilih cat musuh yang paling mengancam area sekutu, dan menara memilih
jenis unit yang paling dibutuhkan berdasarkan komposisi pasukan saat ini.

\subsubsection{Fungsi Kelayakan}

Fungsi kelayakan memastikan bahwa aksi yang dipilih oleh fungsi seleksi
dapat benar-benar dilaksanakan.  Setiap kandidat aksi diperiksa
kelayakannya terlebih dahulu sebelum dieksekusi, misalnya memastikan tile
tujuan dapat dijangkau, bahwa \textit{cooldown} robot sudah siap, serta
bahwa sumber daya yang diperlukan tersedia.

\subsubsection{Fungsi Solusi}

Fungsi solusi dikelola sepenuhnya oleh \textit{game engine}.  Kondisi
solusi terpenuhi apabila salah satu kondisi kemenangan tercapai:
persentase tile bercat tim melebihi 70\%, atau seluruh robot dan menara
lawan berhasil dihancurkan.

\subsubsection{Fungsi Objektif}

Fungsi objektif utama yang ingin dimaksimalkan adalah persentase tile
yang berhasil diwarnai oleh tim:
\begin{equation}
    f_{\text{obj}} =
    \frac{|\text{tile bercat tim sendiri}|}
         {|\text{total tile yang dapat diwarnai}|}
    \times 100\%
    \label{eq:obj}
\end{equation}
Kemenangan diraih apabila $f_{\text{obj}} > 70\%$.  Secara lokal, setiap
robot berusaha memaksimalkan kontribusinya terhadap (\ref{eq:obj}) dengan
selalu memilih aksi yang memberikan peningkatan terbesar terhadap cakupan
wilayah tim pada ronde yang sedang berjalan.

\subsection{Eksplorasi Alternatif Solusi \textit{Greedy}}

\subsubsection{Alternatif 1: Eksplorasi Agresif Tanpa Refill (\texttt{mainbot})}

\paragraph{Heuristik Umum: \textit{Greedy Exploration} dan \textit{Anti-Trailing}}
Secara umum, seluruh unit menggunakan fungsi seleksi pergerakan yang meminimalkan fungsi biaya total $J$ untuk setiap arah $d \in D$:
\begin{equation}
    J(d) = \text{dist}^2(\text{loc} + d, \text{target}) + \text{penalty}_{\text{trail}}(\text{loc} + d)
\end{equation}
Di mana $\text{penalty}_{\text{trail}}$ memberikan bobot penalti sebesar 800 poin untuk tile yang sudah pernah dikunjungi (disimpan dalam memori \textit{trail}). Hal ini mencegah robot terjebak dalam siklus gerakan lokal dan memastikan eksplorasi yang tersebar luas.

\paragraph{Fungsi Seleksi Spesifik Unit \textit{Soldier}}
\textit{Soldier} bertindak sebagai unit pembangun dan pengecat presisi. Fungsi seleksinya dibagi menjadi tiga tahap utama:
\begin{enumerate}
    \item \textbf{Tahap Awal (\textit{Initial Paint}):} Melakukan pembagian tugas (\textit{role}) berdasarkan paritas ID robot ($ID \pmod 2$). Robot memilih tile di sekitar menara induk yang belum berwarna tim untuk menjamin ketersediaan cat di pangkalan.
    \item \textbf{Tahap Konstruksi:} Memprioritaskan kandidat reruntuhan (\textit{ruin}) dengan skor:
    \begin{equation}
        \text{Score}_{\text{ruin}} = 10 \cdot N_{\text{correct\_pattern}} - \frac{\text{dist}^2}{5}
    \end{equation}
    di mana $N_{\text{correct\_pattern}}$ adalah jumlah tile pola yang sudah sesuai.
    \item \textbf{Tahap Eksplorasi Agresif:} Jika cat habis ($\text{paint} \le 10$), fungsi seleksi akan memilih arah yang memaksimalkan jarak dari menara sekutu terdekat (\textit{move away}) untuk melakukan \textit{scouting}.
\end{enumerate}

\paragraph{Fungsi Seleksi Spesifik Unit \textit{Mopper}}
\textit{Mopper} berfokus membersihkan cat musuh tanpa melakukan isi ulang (\textit{refill}). Fungsi seleksi aksi memilih tile serangan berdasarkan kepadatan cat musuh di sekitar tile target:
\begin{equation}
    \text{Score}_{\text{clean}} = (100 - \text{dist}_{\text{base}}) + 5 \cdot N_{\text{adjacent\_enemy\_paint}}
\end{equation}
Untuk pertempuran unit (\textit{combat}), \textit{Mopper} secara \textit{greedy} memilih arah ayunan (\textit{mop swing}) yang menghasilkan hit terbanyak terhadap robot lawan.

\paragraph{Fungsi Seleksi Spesifik Unit \textit{Splasher}}
\textit{Splasher} menggunakan logika \textit{quadrant density} untuk menentukan arah pergerakan. Fungsi seleksi membagi area pandang menjadi empat kuadran dan menghitung densitas cat target (kosong atau musuh):
\begin{equation}
    \text{Density}_q = \frac{\sum \text{Score}_{\text{tile} \in q}}{\text{Count}_{\text{tile} \in q}}
\end{equation}
Robot akan bergerak menuju kuadran dengan $\text{Density}_q$ tertinggi. Dalam pemilihan target serangan, \textit{Splasher} memilih tile yang memaksimalkan dampak area (\textit{splash damage}) dengan bobot 3 untuk cat musuh dan 1 untuk tile kosong.

\paragraph{Fungsi Seleksi Strategis \textit{Tower}}
Menara (\textit{TowerUnit}) mengimplementasikan fungsi seleksi produksi untuk menjaga keseimbangan komposisi pasukan. Keputusan tipe unit yang diproduksi ($T$) didasarkan pada rasio populasi saat ini ($R$):
\begin{itemize}
    \item Jika $R_{\text{mopper}} < 15\%$, pilih \textit{Mopper}.
    \item Jika $R_{\text{soldier}} < 40\%$, pilih \textit{Soldier}.
    \item Jika $R_{\text{splasher}} < 45\%$, pilih \textit{Splasher}.
\end{itemize}
Selain itu, menara melakukan \textit{broadcasting} lokasi secara berkala untuk memfasilitasi koordinasi simetris antar unit, memastikan robot tahu arah mana yang harus dijauhi saat cat mereka habis.

\subsubsection{Alternatif 2: Pendekatan \textit{Macro-Efficiency} dan \textit{Proactive Map Control} (\texttt{altbot1})}

Alternatif kedua ini berfokus pada \textit{map control} secara masif dengan memanfaatkan pergerakan yang aman dan stabil. Pengambilan keputusan \textit{greedy} pada alternatif ini tidak hanya berfokus pada penyerangan musuh, tetapi sangat mementingkan efisiensi cat dan pencegahan kerugian. Untuk mendukung keputusan \textit{greedy} tersebut, bot ini menggunakan perhitungan matematis dan operasi \textit{bitwise} yang kuat. Terdapat tiga heuristik utama yang membedakan alternatif ini:

\paragraph{Heuristik Navigasi Aman dan \textit{Proactive Painting}}

Pada setiap giliran, robot harus memilih satu dari 8 arah mata angin. Fungsi seleksi mengevaluasi setiap petak kandidat menggunakan fungsi objektif \textit{weighted navigation scoring}. Rumus fungsi objektif untuk mengevaluasi arah $d$ adalah sebagai berikut:

\begin{equation}
    Score(d) = W_{arah} + W_{cat} + W_{pusat} - P_{sejarah} - P_{menara}
\end{equation}

Distribusi bobot \textit{Greedy} pada rumus di atas ditentukan sebagai berikut:
\begin{itemize}
    \item $W_{arah} \in \{20, 10, 2\}$: Bobot keselarasan dengan arah makro (\textit{macro direction}). Arah makro ini didapatkan dari prediksi markas musuh menggunakan rumus simetri rotasi pusat:
    \begin{align}
        X_{musuh} &= W - 1 - X_{kita} \\
        Y_{musuh} &= H - 1 - Y_{kita}
    \end{align}
    \textit{(di mana $W$ adalah lebar peta dan $H$ adalah tinggi peta)}
    \item $W_{cat} \in \{8, 10, -15\}$: Pemilihan petak berbasis cat. Berjalan di cat sekutu diberikan bonus kecepatan (+8), tanah kosong didorong untuk ekspansi (+10), dan cat musuh diberi penalti absolut (-15) untuk menghindari terkurasnya persediaan cat.
    \item $W_{pusat} = 5$: Bonus jika langkah tersebut mendekatkan robot ke titik tengah peta.
    \item $P_{sejarah} = 60$: Penalti absolut jika petak sudah ada di memori \textit{Tabu List} lokal untuk mencegah \textit{looping}.
    \item $P_{menara} = 1000$: Penalti mutlak \textit{Tower Starvation}. Bot secara \textit{greedy} menolak petak yang masuk ke dalam radius serang \textit{Defense Tower} musuh untuk mencegah kematian konyol.
\end{itemize}

\paragraph{Heuristik Pembangunan Menara Dinamis Berbasis Rasio}

Ketika sebuah \textit{Soldier} tiba di sebuah puing \textit{ruin}, ia mengambil keputusan \textit{greedy} instan untuk menentukan tipe menara apa yang paling menguntungkan saat itu juga. Untuk menyeimbangkan ekonomi, bot menghitung rasio menara secara global:

\begin{equation}
    Rasio = \frac{\sum PaintTower}{\sum MoneyTower}
\end{equation}

Keputusan \textit{Greedy} mutlak yang diambil adalah jika $Rasio < 2.0$, maka robot \textbf{wajib} membangun \textit{Paint Tower} untuk memastikan \textit{supply chain} pasukan tidak terputus. Jika aman dan rasio terpenuhi, barulah \textit{Money Tower} dibangun.

Selain itu, untuk menghindari kondisi di mana robot berebut membangun satu \textit{Ruin} yang sama, bot memampatkan koordinat 2D menjadi 12-bit integer dan menyimpannya di \textit{SharedArray} dengan masa kedaluwarsa:
\begin{align}
    Val &= ((X + 1) \ll 6) \lor (Y + 1) \\
    Payload &= (Val \ll 16) \lor Ronde_{sekarang} \\
    Usia &= Ronde_{sekarang} - (Payload \land \text{0xFFFF})
\end{align}
Jika $Usia > 20$, maka klaim \textit{greedy} atas puing tersebut dianggap hangus dan dapat diambil alih oleh \textit{Soldier} lain.

\paragraph{Heuristik \textit{Heat-Seeking Splasher} dengan Mekanisme \textit{Anti-Clump}}

Unit \textit{Splasher} dirancang sebagai proyektil penyebar cat. Heuristik serangan areanya mengevaluasi setiap petak 3x3 di sekitarnya. Titik pusat yang memberikan konversi petak kosong atau petak musuh menjadi petak sekutu paling banyak akan dipilih sebagai target tembakan (sebuah pencapaian \textit{local optimum}).

Selain itu, navigasi \textit{Splasher} diberikan bobot negatif tambahan jika bergerak mendekati \textit{Splasher} kawan (\textit{Anti-Clumping}), sehingga mereka secara \textit{greedy} akan menyebar sejauh mungkin satu sama lain bagaikan gas yang mengisi ruang kosong, mengoptimalkan luas area yang diwarnai dalam waktu sesingkat mungkin.

\subsubsection{Alternatif 3: Pendekatan Eksplorasi Simetris Secara Sadar (\texttt{altbot2})}

Alternatif ketiga ini mengedepankan efisiensi informasi melalui pemanfaatan struktur geometri peta. Strategi ini berangkat dari observasi bahwa peta dalam Battlecode 2025 selalu memiliki salah satu dari tiga jenis simetri: rotasional 180$^\circ$, refleksi horizontal, atau refleksi vertikal. Dengan mengasumsikan simetri sejak awal, robot dapat melakukan pencarian lokasi menara musuh secara deterministik tanpa harus menjelajah seluruh peta secara acak.

\paragraph{Heuristik Navigasi dengan Bias Simetri (\textit{Symmetry Bias})}
Pada fase awal (\textit{early game}), setiap robot (khususnya \textit{Soldier}) bertindak sebagai informan. Fungsi seleksi pergerakan memberikan skor prioritas pada arah $d$ yang membawa robot mendekati koordinat "cermin" atau "putar" dari posisi markas sendiri. Skor evaluasi arah dihitung sebagai berikut:
\begin{equation}
    Score_{move}(d) = B_{sym} + B_{paint} - P_{trail} - P_{tower}
\end{equation}
Di mana komponen bobot ditentukan secara \textit{greedy} sebagai berikut:
\begin{itemize}
    \item $B_{sym} \in \{300, 500\}$: Bonus besar jika arah $d$ mendekatkan robot ke kandidat lokasi menara musuh berdasarkan asumsi simetri saat ini.
    \item $B_{paint}$: Bonus tambahan jika ubin di arah tersebut belum diwarnai atau mengandung cat musuh.
    \item $P_{trail} = 800$: Penalti jika ubin sudah ada dalam memori \textit{trail} lokal untuk mendorong eksplorasi ke area baru.
    \item $P_{tower}$: Penalti jika arah tersebut terlalu dekat dengan menara sendiri yang sudah mapan, guna memastikan penyebaran unit.
\end{itemize}

\paragraph{Heuristik Konfirmasi dan Eliminasi Simetri (\textit{Symmetry Breaker})}
Berbeda dengan pendekatan yang mengandalkan memori bersama (\textit{SharedMemory}), alternatif ini menggunakan logika \textit{local-deduction}. Setiap robot menyimpan tiga status kemungkinan simetri: \textit{Rotational}, \textit{Horizontal}, dan \textit{Vertical}. Setiap kali robot memindai fitur peta tertentu (seperti tembok atau reruntuhan), robot melakukan pengecekan validitas:
\begin{enumerate}
    \item Robot menghitung posisi simetris dari fitur tersebut berdasarkan ketiga jenis simetri.
    \item Jika ditemukan bukti yang bertentangan (misalnya, ada tembok di posisi $(x, y)$ namun posisi simetrisnya adalah ruang kosong), maka simetri tersebut diberikan poin diskualifikasi.
    \item Secara \textit{greedy}, robot akan segera mengubah arah pergerakannya menuju kandidat lokasi yang ditunjukkan oleh simetri dengan bukti kesalahan paling sedikit.
\end{enumerate}

\paragraph{Heuristik Prioritas Tindakan Lokal (\textit{Greedy Local Action})}
Dalam setiap ronde, robot mengambil keputusan aksi berdasarkan urutan prioritas skor (tumpukan keputusan \textit{greedy}) tanpa melakukan simulasi jangka panjang untuk menghemat \textit{bytecode}:
\begin{itemize}
    \item \textbf{Prioritas 1 (Skor Maksimal):} Menyelesaikan pola menara (\textit{complete pattern}) jika robot berada di atas \textit{ruin}.
    \item \textbf{Prioritas 2 (Skor Tinggi):} Menyerang atau mewarnai ubin musuh yang berada dalam radius aksi, terutama di dekat \textit{ruin} aktif.
    \item \textbf{Prioritas 3 (Skor Menengah):} Melakukan \textit{Rush} agresif ke arah koordinat menara musuh jika simetri sudah terkonfirmasi melalui observasi lokal.
    \item \textbf{Prioritas 4 (Skor Dasar):} Bergerak secara eksploratif ke arah pusat peta atau area netral untuk memperluas cakupan cat tim.
\end{itemize}

Pendekatan ini sangat efektif pada peta berukuran besar di mana waktu tempuh antar menara menjadi faktor krusial. Dengan mengetahui lokasi lawan lebih awal melalui deduksi simetri, tim dapat melakukan serangan terkoordinasi (\textit{timing attack}) sebelum lawan berhasil membangun ekonomi yang stabil.

\subsection{Analisis Efisiensi dan Efektivitas Alternatif Solusi Greedy}

Bagian ini mengevaluasi kinerja ketiga alternatif strategi \textit{greedy} berdasarkan penggunaan sumber daya komputasi (\textit{bytecode}) dan keberhasilan strategi tersebut dalam berbagai kondisi peta.

\begin{enumerate}
    \item \textbf{Analisis Efisiensi (Penggunaan Sumber Daya)}
    \begin{itemize}
        \item \textbf{Alternatif 1 (\texttt{mainbot}):} Paling hemat energi. Karena logikanya sederhana dan tidak memproses banyak aturan pengisian ulang (\textit{refill}), robot bisa lebih fokus memantau wilayah yang luas. Cara geraknya juga sangat ringan sehingga jarang terjadi kendala teknis. Namun, bot ini berisiko kehilangan banyak pasukan karena robot sering kehabisan cat di area lawan.
        
        \item \textbf{Alternatif 2 (\texttt{altbot1}):} Memiliki beban kerja paling berat. Bot ini menggunakan sistem perhitungan yang sangat detail dan sistem catatan bersama (\textit{SharedArray}) yang kompleks untuk mengatur posisi robot. Meskipun ada sistem pengaman agar bot tidak melampaui batas waktu, bot ini berisiko melambat jika jumlah pasukan sudah sangat banyak.
        
        \item \textbf{Alternatif 3 (\texttt{altbot2}):} Memiliki beban kerja menengah. Bot ini menghabiskan banyak energi di awal permainan untuk membaca pesan dan menebak bentuk peta. Keuntungannya, bot ini cukup hemat dalam penggunaan memori karena setiap unit bekerja secara mandiri setelah mengetahui posisi lawan.
    \end{itemize}

    \item \textbf{Analisis Efektivitas (Pencapaian Tujuan)}
    \begin{itemize}
        \item \textbf{Alternatif 1 (\texttt{mainbot}):} Sangat unggul untuk serangan cepat. Karena robot terus menekan tanpa sering pulang untuk mengisi cat, wilayah lawan bisa dikuasai dengan sangat cepat pada awal permainan.
        
        \item \textbf{Alternatif 2 (\texttt{altbot1}):} Paling tangguh untuk permainan jangka panjang. Bot ini sangat pintar mengatur ekonomi tim melalui perhitungan jumlah menara. Selain itu, robot penyerangnya diatur agar menyebar merata dan tidak menumpuk, sehingga penguasaan wilayah menjadi lebih stabil di seluruh peta.
        
        \item \textbf{Alternatif 3 (\texttt{altbot2}):} Unggul dalam mencari lokasi musuh dan pertempuran taktis. Dengan menebak simetri peta sejak awal, bot ini bisa langsung menuju markas musuh dengan tepat. Dalam pertarungan, bot ini sangat mematikan karena selalu mengincar musuh yang paling mudah dikalahkan terlebih dahulu.
    \end{itemize}
\end{enumerate}

Keberhasilan ketiga strategi ini sangat bergantung pada ukuran peta. Pada peta kecil, strategi serangan cepat seperti Alternatif 1 dan 3 lebih efektif. Namun pada peta besar, sistem kendali wilayah dan pengaturan logistik milik Alternatif 2 jauh lebih bisa diandalkan agar tim tidak kehabisan sumber daya di tengah jalan.

\subsection{Strategi \textit{Greedy} yang Dipilih}

Berdasarkan analisis di atas, strategi yang dipilih sebagai solusi akhir adalah \textbf{Alternatif 1 (\texttt{mainbot})} yang dikombinasikan dengan beberapa fitur unggulan dari \textbf{Alternatif 2 (\texttt{altbot1})}. 

Penggabungan ini dilakukan untuk mengambil kecepatan dan ringannya program dari \texttt{mainbot}, sekaligus memperbaiki masalah logistiknya menggunakan sistem dari \texttt{altbot1}.

\paragraph{Alasan dan Pertimbangan Pemilihan:}
\begin{enumerate}
    \item \textbf{Kecepatan Ekspansi (Momentum):} Dalam kompetisi ini, siapa yang cepat menguasai lahan biasanya akan menang. Kecepatan \texttt{mainbot} dalam menyebar tanpa aturan pengisian cat yang rumit memungkinkannya menguasai wilayah luas sebelum musuh siap bertahan.
    
    \item \textbf{Performa yang Stabil:} Batasan waktu berpikir robot sangat ketat. Logika \texttt{mainbot} yang sederhana memastikan robot tetap bisa bergerak lancar dan menyerang tepat waktu meskipun jumlah pasukan tim sudah sangat banyak. Ini menghindari risiko robot macet yang sering dialami Alternatif 2.
    
    \item \textbf{Ketangguhan di Berbagai Peta:} Strategi ini tidak bergantung pada tebak-tebakan bentuk peta yang berisiko salah. Bot bekerja hanya berdasarkan apa yang dilihat di sekitarnya, sehingga lebih konsisten saat menghadapi bentuk peta yang aneh atau penuh rintangan.
    
    \item \textbf{Penggabungan Keunggulan:} Untuk mengatasi masalah "robot mati perlahan" karena kehabisan cat pada \texttt{mainbot} asli, kami menambahkan sistem \textbf{Manajemen Menara} dari \texttt{altbot1}. Dengan sistem ini, tim akan dipaksa membangun menara penghasil cat saat stok menipis, sehingga pasukan tetap bisa bergerak agresif namun tetap memiliki cadangan logistik yang aman.
\end{enumerate}